\section{Introduction}
\label{sec:intro}

Microscaling (MX) formats, standardised by the Open Compute
Project~\cite{ocp2023mx}, associate a block of $k$ elements with one 8-bit
power-of-two exponent (E8M0) and encode each element in a narrow floating-point
type. Four-bit MX has been shown adequate for a range of inference and even
training workloads~\cite{rouhani2023microscaling}, yet its single shared
power-of-two scale bounds the precision available within a block. Workloads
sensitive to this bound have motivated augmented four-bit formats:
\mtwo{}~\cite{m2xfp2026} allocates per-subgroup metadata to recover precision,
and NVFP4 format replaces the conventional E8M0 block scale with a higher-precision FP8 block scale. These are proposed
and evaluated as data formats; their realisation cost, measured against the formats they would replace on identical hardware, is largely absent from the
literature.

This work examines a further variant that we propose, \fpfourres, in which an activation block
is represented as the sum of two \fpfour{} blocks carrying independent E8M0
scales, while weight blocks remain plain \fpfour. The construction is
asymmetric by design: weight footprint is that of \fpfour, whereas activation
precision approaches that of \fpeight. For inference, where weights dominate
memory and bandwidth, this is a favourable position \emph{if} the arithmetic
can be executed at competitive throughput.

That condition is the subject of the paper. A \fpfourres{} dot product forms
two product terms under two independent scales and therefore needs three
genuine 64-bit vector operands and two independent block scales at once,
exceeding a three-read-port register file. The resulting throughput cannot be
inferred from the format definition and must be measured on synthesised
hardware. We do so by implementing \fpfourres{} together with \fpfour,
\fpeight, and \mtwo{} as independent, format-specialised engines behind one
coprocessor, and
characterising all four on a common core, synthesis flow, and technology node.

\subsection{Related work}
\label{sec:intro-related}

Two RISC-V extensions for MX dot products precede this
work. MXDOTP~\cite{islamoglu2025mxdotp} adds an \fpeight{} dot-product unit to
a Snitch core and supplies block scales through Stream Semantic Registers.
VMXDOTP~\cite{wipfli2026vmxdotp} extends the RISC-V Vector specification to
cover \fpeight{} and \fpfour, resolving operand-port pressure by
time-multiplexing vector register-file reads, and reports
\SI{250}{\giga\flop\per\second} at \SI{1632}{\giga\flop\per\second\per\watt}
for \fpfour{} in \SI{12}{\nano\meter}. This work targets a different point in
the design space: a coprocessor behind the standardised \xif{} interface,
hosted by a compact embedded core, covering formats these extensions do not.
Neither prior extension supports \mtwo{} nor our proposed \fpfourres.

The accumulation datapath common to all three works derives from the anchored
fixed-point scheme of Lutz \emph{et al.}~\cite{lutz2024fp8dot}, which
MXDOTP extended to the two independent block scales of an MX dot product. The
present contribution is the derivation of this scheme for four distinct
formats, each with its own frame geometry.

\mtwo{} was introduced as an algorithm–hardware co-design with a systolic-array
accelerator~\cite{m2xfp2026}; the engine described here is a RISC-V coprocessor
realisation that re-derives the per-subgroup top-1 element within the datapath.

\subsection{Contributions}
\label{sec:intro-contrib}

\begin{enumerate}
  \item \fpfourres, a residue MX format, and its hardware evaluation: on a
        common \SI{45}{\nano\meter} flow, its throughput relative to \fpeight{}
        is \DATAGAP{X}$\times$ (\S\ref{sec:eval-residue}), locating its benefit
        in weight memory (\S\ref{sec:analysis}) rather than compute.
  \item A four-format comparison of \fpfour, \fpeight, \mtwo, and \fpfourres{}
        on one host core and one synthesis flow, removing the process and
        system-scope differences that complicate cross-paper MX comparisons
        (\S\ref{sec:eval}).
  \item A per-format scale-free accumulation frame, with anchor and width
        derived in closed form for each format and shown tight against exact
        rational arithmetic over the full E8M0 scale range
        (\S\ref{sec:arch-frame}).
  \item A coprocessor conforming to the \xif{} interface, portable across
        compliant host cores (\S\ref{sec:bg-xif}).
\end{enumerate}